\documentclass[11pt]{article}
\usepackage[margin=1in]{geometry}
\usepackage{amsmath,amssymb}
\usepackage{enumitem}
\usepackage[T1]{fontenc}

\title{COMP2300 Set 3 Practice Exam Solutions}
\date{}

\begin{document}
\maketitle

\section*{Part A: Multiple Choice}
\begin{enumerate}[label=\textbf{A\arabic*.}]
  \item C
  \item D
  \item B
  \item B
  \item B
  \item C
  \item B
  \item D
  \item B
  \item B
\end{enumerate}

\section*{Part B: Computational and Design Problems}
\begin{enumerate}[label=\textbf{B\arabic*.}]
  \item \textbf{SR Latch}\\
  Using NOR SR latch rules:
  \begin{itemize}[leftmargin=*]
    \item (0,0): hold $Q=0$
    \item (1,0): set $Q=1$
    \item (0,0): hold $Q=1$
    \item (0,1): reset $Q=0$
    \item (0,0): hold $Q=0$
    \item (1,1): invalid (both outputs forced low)
    \item (0,0): indeterminate due to prior invalid state
  \end{itemize}

  \item \textbf{D Latch}\\
  Transparent when CLK=1; hold when CLK=0.
  \begin{center}
  \begin{tabular}{c|ccccc}
  Interval & 1 & 2 & 3 & 4 & 5 \\ \hline
  CLK & 0 & 1 & 1 & 0 & 1 \\
  D   & 1 & 0 & 1 & 1 & 0 \\
  Q   & 0 & 0 & 1 & 1 & 0 \\
  \end{tabular}
  \end{center}

  \item \textbf{D Flip-Flop}\\
  Samples on rising edges. Starting $Q=1$:
  \[
  Q(t_1)=0,\ Q(t_2)=1,\ Q(t_3)=1,\ Q(t_4)=0
  \]

  \item \textbf{Registers and Memory}\\
  \begin{enumerate}[label=(\alph*)]
    \item $\log_2(64\text{K})=16$ address bits.
    \item Capacity $= 64\text{K} \times 16$ bits $= 1,048,576$ bits $= 131,072$ bytes (128 KiB).
    \item Byte-addressable capacity $= 64\text{K} \times 8$ bits $= 524,288$ bits $= 65,536$ bytes (64 KiB).
  \end{enumerate}

  \item \textbf{Timing constraints}\\
  \begin{enumerate}[label=(\alph*)]
    \item $T_c = t_{pcq} + t_{pd} + t_{setup} = 0.4 + 6.0 + 0.2 = 6.6$ ns.
    \item $f_{max} = 1/T_c \approx 1/6.6\text{ ns} = 151.5$ MHz.
    \item Aperture time $= t_{setup}+t_{hold} = 0.2 + 0.1 = 0.3$ ns.
  \end{enumerate}

  \item \textbf{Pipeline throughput}\\
  \begin{enumerate}[label=(\alph*)]
    \item $T_c = 0.3 + \max(4,3,5) + 0.2 = 5.5$ ns.
    \item Cycles $= K + (N-1) = 3 + 19 = 22$. Time $= 22 \times 5.5$ ns $= 121$ ns.
  \end{enumerate}
\end{enumerate}

\section*{Part C: Past Exam Questions}
\begin{enumerate}[label=\textbf{C\arabic*.}]
  \item \textbf{Moore FSM concept}\\
  A Moore FSM is a finite state machine whose outputs depend only on the current state, not directly on inputs. Example: a traffic light controller where each state encodes the light color.

  \item \textbf{FSM design (three consecutive 1s)}\\
  \begin{enumerate}[label=(\alph*)]
    \item States represent count of consecutive 1s: $S_0$ (0), $S_1$ (1), $S_2$ (2), $S_3$ (3 or more). Output $X=1$ only in $S_3$.
    Transitions: on $A=1$, $S_0\to S_1\to S_2\to S_3\to S_3$; on $A=0$, any state $\to S_0$.
    \item Example encoding: $S_0=00$, $S_1=01$, $S_2=10$, $S_3=11$.
    \item Next-state equations (one valid set):
    \[
    D_1 = A \cdot (Q_1 + Q_0),\quad D_0 = A \cdot (Q_1 + \overline{Q_0})
    \]
    Output: $X=Q_1 Q_0$.
    \item Schematic: input $A$ into next-state logic; state register is two D flip-flops; outputs from register feed output logic for $X$.
  \end{enumerate}

  \item \textbf{Pipeline timing}\\
  \begin{enumerate}[label=(\alph*)]
    \item $T_c = 0.2 + \max(3.4,3.3,3.7) + 0.1 = 4.0$ ns.
    \item $K=3$, $N=12$ cycles $= 3 + 11 = 14$. Time $= 14 \times 4.0$ ns $= 56$ ns.
    \item $N=10^9$ cycles $= 10^9 + 2$. Time $\approx 4.0\text{ ns} \times (10^9+2) \approx 4.000000008$ s.
    \item Four pipelines: $2.5\times 10^8$ tokens each. Time $\approx 4.0\text{ ns} \times (2.5\times 10^8 + 2) \approx 1.000000008$ s.
  \end{enumerate}
\end{enumerate}

\end{document}
